\section{Engenharia Elétrica no cenário de São Luís e do Maranhão}\label{sec:cenario}

A modernização dos laboratórios do DEE está diretamente ligada ao contexto econômico e tecnológico do Maranhão. Os indicadores a seguir mostram por que a Engenharia Elétrica é estratégica para a formação profissional e para projetos de pesquisa aplicada e P\&D no estado.

\begin{center}
\begin{tabularx}{\textwidth}{>{\centering\arraybackslash}X >{\centering\arraybackslash}X >{\centering\arraybackslash}X}
\begin{tcolorbox}[colback=cinzaclaro,colframe=azulpetroleo!40,boxrule=0.5pt,arc=1.2mm,left=2mm,right=2mm,top=2mm,bottom=2mm]
{\sffamily\fontsize{22}{24}\selectfont\bfseries\color{azulpetroleo} 8,8\%}\\[-0.5mm]
{\small\bfseries crescimento da indústria maranhense}\\
{\footnotesize 3º trimestre de 2025}
\end{tcolorbox}
&
\begin{tcolorbox}[colback=cinzaclaro,colframe=azulpetroleo!40,boxrule=0.5pt,arc=1.2mm,left=2mm,right=2mm,top=2mm,bottom=2mm]
{\sffamily\fontsize{22}{24}\selectfont\bfseries\color{azulpetroleo} 22,1\%}\\[-0.5mm]
{\small\bfseries avanço da indústria de transformação}\\
{\footnotesize mesmo período}
\end{tcolorbox}
&
\begin{tcolorbox}[colback=cinzaclaro,colframe=azulpetroleo!40,boxrule=0.5pt,arc=1.2mm,left=2mm,right=2mm,top=2mm,bottom=2mm]
{\sffamily\fontsize{22}{24}\selectfont\bfseries\color{azulpetroleo} 325.778}\\[-0.5mm]
{\small\bfseries empregos formais em São Luís}\\
{\footnotesize junho de 2025}
\end{tcolorbox}
\\
\begin{tcolorbox}[colback=cinzaclaro,colframe=azulpetroleo!40,boxrule=0.5pt,arc=1.2mm,left=2mm,right=2mm,top=2mm,bottom=2mm]
{\sffamily\fontsize{22}{24}\selectfont\bfseries\color{azulpetroleo} 51.513}\\[-0.5mm]
{\small\bfseries empregos industriais em São Luís}\\
{\footnotesize cerca de 47\% do total do estado}
\end{tcolorbox}
&
\begin{tcolorbox}[colback=cinzaclaro,colframe=azulpetroleo!40,boxrule=0.5pt,arc=1.2mm,left=2mm,right=2mm,top=2mm,bottom=2mm]
{\sffamily\fontsize{22}{24}\selectfont\bfseries\color{azulpetroleo} 36,84 mi t}\\[-0.5mm]
{\small\bfseries movimentação do Porto do Itaqui}\\
{\footnotesize recorde em 2025}
\end{tcolorbox}
&
\begin{tcolorbox}[colback=cinzaclaro,colframe=azulpetroleo!40,boxrule=0.5pt,arc=1.2mm,left=2mm,right=2mm,top=2mm,bottom=2mm]
{\sffamily\fontsize{22}{24}\selectfont\bfseries\color{azulpetroleo} 1.427 MW}\\[-0.5mm]
{\small\bfseries capacidade instalada do Complexo Parnaíba}\\
{\footnotesize relevância energética do estado}
\end{tcolorbox}
\end{tabularx}
\end{center}

\begin{destaque}
\textbf{Expansão da infraestrutura elétrica.} O Projeto Tangará prevê \textbf{354 km} de linhas de transmissão, \textbf{2.300 MVA} adicionais de capacidade de transformação e investimento estimado em \textbf{R\$ 1,1 bilhão}. Somam-se a isso os investimentos associados ao sistema Graça Aranha--Silvânia, voltado ao escoamento de geração renovável do Norte e Nordeste.
\end{destaque}

Esse contexto reforça a demanda por profissionais capazes de projetar, operar, proteger, supervisionar, ensaiar, medir e modernizar instalações, máquinas, acionamentos, processos automatizados, sistemas de potência, redes industriais e infraestrutura energética. Para o IFMA, isso significa formar profissionais alinhados às necessidades reais do Maranhão e ampliar a possibilidade de participação do DEE em projetos de pesquisa aplicada e desenvolvimento tecnológico.

{\footnotesize\color{cinzamedio}
\textbf{Fontes dos dados:} Instituto Maranhense de Estudos Socioeconômicos e Cartográficos (IMESC); Observatório da Indústria/FIEMA--CNI; Empresa Maranhense de Administração Portuária (EMAP); Ministério de Minas e Energia (MME); Secretaria de Estado de Indústria e Comércio do Maranhão (SEINC). Dados consultados em 2026.}
